\section{Shortening all comparison coordinates}
\label{sec:shortening}

Return to the flat bundle
\[
 Z=(\widetilde M\times\widehat A)/\Gamma\longrightarrow A,
 \qquad q=n+\dim A,
\]
where $M^n$ and $A$ are closed and $\widehat A\to A$ is a principal left
$\Gamma$-cover. The covering projection to $M\times A$ induces the splitting
$TZ=E_M\oplus E_A$ described in \cref{thm:comparison-data}. Here $E_M$ is
tangent to the $\widetilde M$ fibres, while $E_A$ is the flat lift of the
parameter directions. Local choices of sheets over $A$ identify the fibres
with $\widetilde M$; the transition maps are fixed deck transformations.
We first change the metric in the fibre directions and add circles. The
metric on $E_A$ is introduced only in the final assembly step.

A metric with $b$ stabilizing circles has, locally over $A$, the form
\begin{equation}\label{eq:fibre-stabilized-metric}
 g_h(x,a)+\sum_{\nu=1}^b m_\nu(x,a)^2d\theta_\nu^2.
\end{equation}
Here $g_h$ acts on the $n$ fibre directions, the positive warps depend only on $x,a$, and the metric is invariant under all circle translations. At every fixed finite stage, the coefficients and all mixed derivatives will be bounded, and the metric will be uniformly equivalent to the lifted reference product. These bounds may depend on the stage.

Until the parameter metric is assembled in
\cref{lem:parameter-scaling}, every differential, gradient and scalar
curvature is computed on the fibre $\widetilde M\times\T^b$ with
$a\in A$ fixed. The functions and metrics may vary smoothly with $a$;
this variation is measured separately by the mixed derivative bounds.
Thus a torus-invariant fibre differential lies in the original $n$
fibre directions, even though its full differential on $Z$ may
have a nonzero parameter component.

\subsection{The cost of one pass}

A pass processes each of $q$ comparison functions once. In this section,
$L$ denotes a bound for their current fibre differentials. The scalar
estimate for a pass depends on $n$ and $\log q$, even though the full fibre
dimension includes all the circles already introduced.

\begin{lemma}[One pass]\label{lem:one-pass}
Let $q\ge2$ and $L,w>0$. On a uniformly controlled metric of the form
\eqref{eq:fibre-stabilized-metric}, suppose smooth bounded torus-invariant
functions $\ell_1,\ldots,\ell_q$ have bounded derivatives of every order.
For a smooth family over a compact parameter manifold, assume in each of
finitely many parameter charts common bounds for every mixed fibre and
parameter derivative of the metric coefficients, warps, and functions,
and a common equivalence to the reference fibre product. Suppose
\[
 \abs{d\ell_i}_{g_*}\le L,
\]
where $g_*$ denotes the metric on $E_M$ at the start of the pass. Then one
can add $q$ circles and replace these functions by $\ell_i^+$ so that
\begin{equation}\label{eq:pass-function-bounds}
 \abs{\ell_i^+-\ell_i}<\frac\pi2w,
 \qquad \abs{d\ell_i^+}_{g_h^+}\le\frac L4,
 \qquad g_h^+\ge g_*.
\end{equation}
For an absolute constant $C_0$, the full fibre scalar curvature satisfies
\begin{equation}\label{eq:pass-scalar-bound}
 \Scal_{g^+}\ge\Scal_g
       -C_0(n+1)(2+\log q)\frac{L^2}{w^2}.
\end{equation}
All data at the new finite stage remain uniformly controlled. For a compact
parameter family satisfying the stated hypotheses, the new mixed derivatives
are also uniformly bounded.
\end{lemma}

\begin{proof}
Set
\begin{equation}\label{eq:pass-parameters}
 \eps=L/4,\qquad a_0=(2+\log q)^{-1},\qquad
 \beta=1-a_0,\qquad s=q^{-3},\qquad B_0=\frac{4L}{a_0w}.
\end{equation}
Choose a smooth nondecreasing cutoff $\phi$ which is zero for $t\le1/4$, one for $t\ge1/2$, and satisfies $0\le\phi'\le10$. Let
\begin{equation}\label{eq:pass-cutoff}
 N(t)^2=s^2+(1-s^2)\phi(t).
\end{equation}
Writing $\mathfrak p=tN'/N$, we have
\begin{equation}\label{eq:pass-cutoff-positivity}
 \mathfrak p\ge0,\qquad
 \mathfrak p N(t)^2t^2
       =\tfrac12 t^3(1-s^2)\phi'(t)\le5/8<1.
\end{equation}
The expression vanishes outside $[1/4,1/2]$, which justifies the displayed bound.

For one coordinate, denote the original function by $\ell$ and the current full fibre metric by $g$. Solve \eqref{eq:graph-equation} with these choices and replace
\[
 g\longmapsto g^+=g+N(t)^2df^2+v^2d\theta^2,
 \qquad \ell\longmapsto\ell^+=\eps f.
\]
\Cref{thm:uniform-graph} applies on $\widetilde M\times\T^b$. Its
uniqueness gives torus invariance and compatibility with the parameter
trivializations. Consequently $df$ vanishes on the circle directions, so
the metric retains the form \eqref{eq:fibre-stabilized-metric}. Its block
on $E_M$ has only increased. Since
\[
 \abs{df}_{g^+}=vt\le1
\]
(the inequality is immediate for $t<1/2$, and $N=1$ for $t\ge1/2$), the new coordinate has derivative norm at most $L/4$. Subsequent steps can only improve this bound. The tangent band gives the height estimate in \eqref{eq:pass-function-bounds}.

It remains to estimate the scalar loss in a way that can be summed. At $t>0$, put $z=e_1r$. Differentiating $\lambda(r)=B_0(r^2+s^2)^{\beta/2}$ along $U=re_1$ gives
\[
 \frac{U(\lambda)}{\lambda}=\gamma z,
 \qquad \gamma=\frac{\beta r^2}{r^2+s^2}\in[0,\beta].
\]
The part of $-2U(D)$ contributed by this derivative is $-2\gamma zD$. Hence
\begin{equation}\label{eq:pass-square-absorption}
 D^2+z^2-2\gamma zD
       \ge(1-\gamma^2)D^2\ge a_0D^2.
\end{equation}
By \eqref{eq:pass-cutoff-positivity}, the final term of \eqref{eq:scalar-exact} is nonnegative.

Measure the vector $U$ in the metric at the start of the pass, and set
\[
 b_U=\abs{U}_{g_*}.
\]
Since the current metric on $E_M$ is at least $g_*$, one has $b_U\le r$.
The input coordinate still satisfies $\abs{d\ell}_{g_*}\le L$, so
$U(\ell)\le Lb_U$, while $U(f)=rt$. Differentiating the tangent in the
equation therefore contributes an adverse term bounded above by
\begin{equation}\label{eq:pass-adverse-term}
 \frac2w\left(\lambda+\frac{D^2}{\lambda}\right)
                      (Lb_U-\eps rt)_+.
\end{equation}
Since $0\le r<1$ and $\beta<1$,
\[
 \lambda=B_0(r^2+s^2)^{\beta/2}\ge B_0r.
\]
Thus the $D^2$ part of \eqref{eq:pass-adverse-term} is at most
\[
 \frac{2L}{wB_0}D^2=\frac{a_0}{2}D^2,
\]
which is absorbed by \eqref{eq:pass-square-absorption}.

For the remaining part, first suppose $r<s$. Then
\begin{equation}\label{eq:pass-small-slope}
 \lambda b_U\le \sqrt2 B_0s^{\beta+1}\le\frac{\sqrt2 B_0}{q}.
\end{equation}
Here $s=q^{-3}$ and $\beta>0$. If instead $r\ge s$, then $t\ge1/4$: for $t<1/4$ one has $N=s$ and $r< s/4$. On the adverse region where $Lb_U-\eps rt>0$, it follows that
\[
 b_U>\frac{rt}{4}\ge\frac r{16}.
\]
Moreover,
\[
 r^{-a_0}\le s^{-a_0}
       =\exp\left(\frac{3\log q}{2+\log q}\right)\le e^3.
\]
Consequently, on that region,
\begin{equation}\label{eq:pass-large-slope}
 \lambda b_U\le\sqrt2 B_0r^{1-a_0}b_U
            \le16\sqrt2 e^3B_0b_U^2.
\end{equation}
Combining these estimates and substituting $B_0$ from \eqref{eq:pass-parameters} gives the one-coordinate inequality
\begin{equation}\label{eq:coordinate-scalar-cost}
 \Scal_{g^+}\ge\Scal_g
      -C_0\frac{L^2}{a_0w^2}\left(\frac1q+b_U^2\right).
\end{equation}
For example $C_0=128\sqrt2 e^3$ is sufficient. The inequality extends by continuity to the boundary of $\{t>0\}$. On the interior of $\{t=0\}$ the graph is constant and the update adds a constant circle, so it holds there as well.

Finally sum \eqref{eq:coordinate-scalar-cost} over the $q$ coordinates.
At a fixed point of $Z$, the inverse metrics of the $E_M$ blocks all act
on the same $n$-dimensional space $E_M^*$. If $U_i$ is the vector used at the $i$th
step, \eqref{eq:graph-inverse-volume} gives the telescoping identity
\begin{equation}\label{eq:pass-trace-telescope}
 \sum_{i=1}^q U_i\otimes U_i
       =g_*^{-1}-(g_h^+)^{-1}\le g_*^{-1}.
\end{equation}
Taking its $g_*$-trace gives $\sum_i b_{U_i}^2\le n$. The $q$ terms $1/q$ sum to $1$, proving \eqref{eq:pass-scalar-bound}. In particular the growing number of old circle directions does not appear in this trace. Uniform control and smooth parameter dependence at the new stage follow from \cref{thm:uniform-graph} applied successively a finite number of times.
\end{proof}

\subsection{Finite iteration and the parameter directions}

We shall use the following elementary assembly observation.

\begin{lemma}[Magnifying the parameter metric]\label{lem:parameter-scaling}
Fix a finite-stage smooth family of metrics of the form
\eqref{eq:fibre-stabilized-metric}, equivariant under the flat transition maps over the closed manifold $A$, with all mixed derivatives bounded and uniform equivalence to a reference product. For any fixed metric $g_A$ on $A$, set
\[
 g_P=g_h+P^2g_A+\sum_{\nu=1}^b m_\nu^2d\theta_\nu^2
 \quad\text{on }Z\times\T^b,
\]
using the flat splitting and no mixed parameter terms. As $P\to\infty$, $\Scal_{g_P}$ converges uniformly to the scalar curvature of the fibre metrics. If $H:Z\to\R^q$ has bounded mixed derivatives, its differential in the parameter directions has norm tending uniformly to zero. Each $g_P$ is complete and uniformly controlled.
\end{lemma}

\begin{proof}
In a parameter chart about $a_*$ write $a=a_*+P^{-1}y$. The parameter block becomes $g_A(a_*+P^{-1}y)$, and the other metric coefficients are the corresponding fibre coefficients evaluated at $a_*+P^{-1}y$. All parameter derivatives through order two tend uniformly to zero. Their limits, at $y=0$, are the two-jets of the fibre metric with $a=a_*$ frozen, times the constant Euclidean parameter block $g_A(a_*)$. Inverses remain uniformly bounded for this fixed stage. The scalar-curvature formula is a continuous expression in these metric two-jets and their inverses, so its limit is uniformly the fibre scalar curvature. The finitely many parameter charts and the assumed bounds make the convergence uniform over all of $Z$.

The parameter differential of $H$ acquires the factor $P^{-1}$, proving the second assertion. The reference product on $Z$ is the pullback of a metric on the closed manifold $M\times A$ and is complete. The assembled metric, with its finite number of circle blocks, is uniformly equivalent to this complete reference product and has bounded coefficient derivatives.
\end{proof}

\begin{theorem}[Logarithmic shortening]\label{thm:shortening}
Suppose $\Scal_M\ge\kappa>0$. On the above bundle $Z$, let
$G:Z\to\R^q$ be smooth and bounded with all mixed derivatives bounded, and suppose $L_0>0$ and
\[
 \sup_{1\le i\le q}\abs{d_{\widetilde M}G_i}\le L_0,
 \qquad q\ge2.
\]
There is a constant $C_n$ depending only on $n$ such that, if
\begin{equation}\label{eq:shortening-smallness}
 C_n(2+\log q)L_0^2\le\kappa/4,
\end{equation}
then for every finite integer $k\ge1$ there are a complete stabilized metric
\[
 \mathfrak g_k=h_k+\sum_{\nu=1}^{kq}m_{k,\nu}^2d\theta_\nu^2
 \quad\text{on }Z\times\T^{kq}
\]
and a smooth bounded map $G^{(k)}:Z\to\R^q$ satisfying
\begin{equation}\label{eq:shortening-output}
 \Scal_{\mathfrak g_k}\ge\kappa/2,\qquad
 \norm{G^{(k)}-G}_{\infty}<\pi/10,\qquad
 \abs{dG^{(k)}}_{\rm op}\le\sqrt{q+1}\,4^{-k}L_0.
\end{equation}
The infinity norm in the middle inequality is the target coordinate maximum, followed by the supremum over $Z$. All metric coefficients, warps, and maps have uniform bounds of every order at each fixed finite stage; their equivalence and derivative constants may depend on that stage.
\end{theorem}

\begin{proof}
Start with the lifted metric of $M$ and no circles, and perform passes indexed by $j=0,\ldots,k-1$ with
\begin{equation}\label{eq:pass-schedule}
 L_j=4^{-j}L_0,\qquad w_j=2^{-j}/10.
\end{equation}
\Cref{lem:one-pass} applies inductively because every coordinate derivative is divided by at least four in each pass. Every stage is finite, so the uniform control required for its successor follows from the last assertion of that Lemma. All local constructions match under the flat parameter transitions by uniqueness of the graph solution.

For any number of passes, the total scalar loss is at most
\[
 C_0(n+1)(2+\log q)\sum_{j\ge0}\frac{L_j^2}{w_j^2}
   =\frac{400}{3}C_0(n+1)(2+\log q)L_0^2.
\]
Choose $C_n=(400/3)C_0(n+1)$. Under
\eqref{eq:shortening-smallness}, all the fibre scalar curvatures are at least $3\kappa/4$. Each coordinate is modified exactly once per pass, so
\[
 \norm{G^{(k)}-G}_{\infty}
   <\frac\pi2\sum_{j=0}^{k-1}w_j<\frac\pi{10}.
\]
No coordinate count enters this estimate. The fibrewise Euclidean operator norm is at most $\sqrt q\,L_k$, since each of the $q$ rows has norm at most $L_k=4^{-k}L_0$.

At this fixed stage apply \cref{lem:parameter-scaling}, making the parameter scale so large that the scalar-curvature loss from assembly is at most $\kappa/4$ and the parameter part of $dG^{(k)}$ has operator norm at most $L_k$. The domain's fibre and parameter blocks are orthogonal. Cauchy--Schwarz therefore bounds the full differential by
$\sqrt{qL_k^2+L_k^2}$, proving \eqref{eq:shortening-output}. The resulting base metric on $Z$ is $h_k$; all circle warps remain positive and independent of circle variables. The completeness and finite-stage bounds follow from the same assembly Lemma.
\end{proof}

\subsection{A compact degree region}

\begin{corollary}[Stabilized contraction at a fixed scale]
\label{cor:stabilized-contraction}
Assume $\Scal_M\ge\kappa>0$ and take the comparison data of
\cref{thm:comparison-data} at scale $R$, so that
\[
 L_0=J/R,\qquad \log q\le JR.
\]
For all sufficiently large $R$, fix these data once and for all. Then for every finite $k\ge1$ there is a stabilized metric $\mathfrak g_k$ on $Z\times\T^{kq}$ with scalar curvature at least $\kappa/2$, and a smooth map
\[
 F_k:Z\longrightarrow S^q
\]
equal to a single value outside a compact set, with nonzero degree, such that
\begin{equation}\label{eq:stabilized-contraction-bound}
 \abs{dF_k}_{\rm op}\le C(q)4^{-k}J/R.
\end{equation}
The sphere has its standard unit metric. The constant $C(q)$ is independent of $k$. If desired, the circle lengths may be increased so that the projection to the standard $\T^{kq}$ also has operator norm at most $4^{-k}J/R$; this does not change scalar curvature or the bound on $dF_k$.
\end{corollary}

\begin{proof}
The left-hand side of \eqref{eq:shortening-smallness} is at most
\[
 C_n(2+JR)J^2/R^2\longrightarrow0.
\]
Choose a sufficiently large finite $R$ once. In particular $q$, $A$, and the original comparison map are now fixed while $k$ remains arbitrary. Apply \cref{thm:shortening} at this fixed scale.

Choose a smooth map $Q_q:\R^q\to S^q$ which is constant outside the Euclidean unit ball and has degree one on the one-point compactification. It can be chosen with the north pole as a regular value having the sole preimage $0$. For example, in polar coordinates use
\[
 y\longmapsto
 \left(\sin\rho(\abs y)\frac{y}{\abs y},\cos\rho(\abs y)\right),
\]
where $\rho(r)$ is a positive linear function near zero, increases to $\pi$, and is smoothly constant at $\pi$ in a collar of $r=1$ and beyond; choose the orientation, or compose with a reflection, to give degree one. The expression is smooth at zero because $\rho$ is linear there. Let $D_q<\infty$ bound its differential.

By the comparison theorem there is a precompact open $\Omega\subset Z$ with
$\deg(G,\Omega,0)\ne0$ and $\norm{G}_{\infty}>3$ on $\partial\Omega$. The homotopy
\[
 G+u(G^{(k)}-G),\qquad 0\le u\le1,
\]
stays outside the unit ball on $\partial\Omega$, by
\eqref{eq:shortening-output}. Define
\[
 F_k=Q_q\circ G^{(k)}\quad\text{on }\Omega,
 \qquad F_k=Q_q(\infty)\quad\text{on }Z\setminus\Omega.
\]
The boundary is compact and its image is separated from the unit ball. Thus $Q_q\circ G^{(k)}$ is constant on a neighborhood of the boundary, and this definition is smooth even if $\partial\Omega$ is not smooth. The map is constant outside the compact set $\overline\Omega$.

Homotopy invariance of local degree gives
\[
 \deg(G^{(k)},\Omega,0)=\deg(G,\Omega,0)\ne0.
\]
The degree of $F_k$ is this same integer: count preimages of a regular target sufficiently near the north pole, where $Q_q$ is an orientation-preserving local diffeomorphism, or equivalently compose the relative degree classes. Hence $F_k$ has nonzero degree. Its differential is bounded by
$D_q\sqrt{q+1}\,4^{-k}J/R$, establishing
\eqref{eq:stabilized-contraction-bound} with $C(q)=D_q\sqrt{q+1}$.

For the last assertion, multiply each circle warp by a sufficiently large positive constant. Because all coefficients are independent of the circle variables and the circle blocks are mutually orthogonal, this is locally a constant change of each circle coordinate. It preserves scalar curvature. The base metric and the norm of $dF_k$ are unchanged. At the fixed stage all warps have positive uniform lower bounds, so their constant multipliers can make every inverse circle length at most $4^{-k}J/R$. The operator norm of the projection to the product flat torus is the maximum of these inverse lengths, with no factor depending on the number of circles.
\end{proof}

We have now constructed maps of nonzero degree from $Z$ to $S^q$.
Their differentials tend to zero as the finite pass count grows, with
$q$ fixed. The scalar lower bound holds on $Z\times\T^{kq}$. The
remaining obstruction must therefore be independent of the number of
stabilizing circles. We establish it in \cref{sec:obstruction}.
